%═════════════════════════════════════════════════════════
\section{Introduction}\label{sec:Introduction}
%═════════════════════════════════════════════════════════

Accurate volatility forecasting is central to modern finance, underpinning portfolio allocation, derivative pricing, and risk management. The relevant object is not only the conditional variance of returns but the full tail distribution, since risk measures such as Value-at-Risk (VaR) depend directly on the probability of extreme losses. Over the past four decades, a large econometric literature has developed around volatility models that exploit well-documented empirical regularities, including volatility clustering, persistence, and heterogeneous dynamics across time horizons. Seminal contributions include the ARCH model of \citet{engle1982autoregressive}, which formalized conditional heteroskedasticity, and the GARCH model of \citet{Bollerslev1986}, which captures persistent volatility dynamics through a parsimonious recursion. Later extensions incorporate asymmetric responses to shocks \citep{Glosten1993} and long-memory behavior in realized volatility \citep{corsi2009simple}. Heavy-tailed innovation distributions such as the Student-t or skew-t \citep{bollerslev1987conditionally, Hansen1994} are commonly adopted to model return distributions. Together, these approaches encode substantial domain knowledge and remain among the most widely used tools for volatility and risk forecasting.

Despite their empirical success, even well-established models remain approximations of a complex and evolving financial system. In practice, the true data-generating process may contain dynamics that are unknown or too complex to represent within a tractable parametric specification \citep{poon2003forecasting, Rossi2021}. Moreover, financial markets are inherently non-stationary, meaning that the statistical properties of returns evolve over time \citep{StockWatson1996, Rossi2013}. Parameters governing volatility persistence, scale, and distributional shape may drift gradually with changes in the environment. In essence, non-stationarity reflects the presence of structural breaks and regime changes whose timing and character are unknown in advance \citep{clements2011forecastingbreaks, Rossi2013}. Consequently, even empirically successful models that capture the dominant stylized facts often leave systematic patterns in their forecast errors, indicating residual dynamics not captured by the underlying parametric structure \citep{poon2003forecasting}. While these remaining dynamics may contain useful predictive information, their evolution is typically unknown and difficult to specify parametrically.

This raises the central question of this paper:

Can variance and tail-risk forecasts be improved by learning from forecast errors, without imposing a specific parametric evolution mechanism?

A natural framework for this challenge is online learning theory (OLT), a branch of statistical learning and optimization explicitly designed for settings where no probabilistic assumptions are made about the data-generating process. Forecasting problems are formulated as convex online optimization tasks, where predictions are updated sequentially to minimize a convex loss function as new observations arrive \citep{cesa2006prediction, ShalevShwartz2012, hazan2016introduction}. A central result is that such algorithms achieve sublinear regret, meaning average forecasting performance converges to that of the best fixed predictor chosen in hindsight. These guarantees hold even under non-stationarity or unknown complex dynamics. Yet, OLT has received limited attention in the volatility forecasting literature, motivating a framework that exploits forecast error information from volatility models using online learning techniques to improve volatility and tail-risk forecasts.

We therefore propose \textit{Model-Assisted Online Learning} (MAOL), a framework that combines classical econometric modeling with OLT. The central idea is to treat an established volatility model as a structural filter that captures the dominant and well-understood components of volatility dynamics documented in the literature; any remaining predictive information then must reside in the sequence of forecast errors. For these residual dynamics, outside the scope of the underlying parametric specification, the MAOL framework applies online learning techniques to adapt forecasts sequentially based on observed forecast errors. Similarly, for distributional forecasts, MAOL starts from a chosen return distribution and sequentially recalibrates it.

The framework has three layers, each driven by the parameter-free Continuous Coin Betting (COCOB) optimizer \citep{orabona2017training}, which avoids manual learning-rate tuning. Layer 1 updates the parameters of the base volatility model observation by observation using the gradient of the QLIKE loss, allowing the model to continuously adapt to structural changes in volatility dynamics without requiring repeated re-estimation on rolling or expanding samples. Layer 2 applies a multiplicative bias correction that learns a variance-level multiplier online to correct slow-moving shifts in unconditional volatility that the recursion alone cannot absorb. Separating level from shape allows online updates to target each component with the appropriate feedback, preserving the interpretability of the dynamic parameters. Layer 3 recalibrates the predictive distribution in probability integral transform (PIT) space via online gradient descent on the pinball loss. Deviations from uniformity in PIT values provide evidence of distributional misspecification \citep{diebold1998evaluating, Berkowitz2001}. MAOL uses this information to sequentially correct the predictive distribution and learn Value-at-Risk (VaR) quantiles.

The literature has addressed model instability through several approaches. Markov regime-switching and smooth-transition GARCH models allow parameters to take different values across states \citep{Hamilton1989, CrealKoopmanLucas2013, harvey2013dynamic}, but require specifying the number of regimes and the transition law. Nonparametric and machine-learning methods such as neural networks avoid imposing a functional form and have been shown to improve predictive accuracy \citep{Buhlmann2002, Donaldson1997}, but abstract from the stylized facts embedded in classical volatility models. Within the parametric tradition, hybrid approaches such as LASSO-regularized HAR and realized-measure combinations improve forecasting performance \citep{AudrinoKnaus2016, BezerraAlbuquerque2017} while retaining interpretability. In contrast, MAOL requires no tuning, preserves the underlying volatility model's economic interpretation, and adapts sequentially without committing to a specific parametric evolution mechanism.

Among existing approaches, the volatility forecasting literature most closely related to ours are \citet{ZhuoMorimoto2024} and \citet{TapiaKristjanpoller2022}, which apply support vector regression to HAR residuals and long short-term memory (LSTM) networks to volatility model residuals, respectively, finding that forecast errors contain exploitable nonlinear information, thereby improving forecasts. However, both hybrid approaches still rely on batch estimation and extensive design choices without theoretical guarantees.

Sequential methods such as rolling windows, recursive quasi-maximum likelihood (QML), and adaptive volatility estimation update parameters as new observations arrive \citep{PesaranTimmermann2007, Dahlhaus2007, Werge2022}, but remain focused on parameter estimation within a fixed parametric specification and are computationally more expensive.

Existing approaches thus face a fundamental tension between interpretability, flexibility, computational efficiency, and sequential adaptivity, a gap that online learning is uniquely positioned to fill. Applications of OLT to financial time series remain limited, with existing work focusing on ARMA parameter updating \citep{AnavaHazanMannorShamir2013, liu2016onlinearima}, portfolio allocation \citep{stoltz2005internal, li2013onlineportfolio}, and covariance estimation \citep{johansson2024covariance}, leaving volatility and tail-risk forecasting largely unexplored.

The proposed MAOL framework bridges these strands of literature by combining the structural interpretability of econometric volatility models with the sequential adaptivity of online convex optimization. Like the hybrid approaches of \citet{ZhuoMorimoto2024} and \citet{TapiaKristjanpoller2022}, it exploits residual predictive information left by the base model. However, it does so sequentially, and with formal regret guarantees, which is a stronger performance criterion than the consistency or asymptotic efficiency results that anchor the recursive estimation literature.

A related literature improves predictive distributions using forecast errors through conformal prediction and sequential calibration methods \citep{FosterVohra1998, kuleshov2018calibrated, Vovk2020, GibbsCandes2021, GibbsCandes2024, Deshpande2024}. Unlike these approaches, which recalibrate predictions around a fixed model, MAOL sequentially updates both the volatility model and the predictive distribution.

MAOL employs COCOB \citep{orabona2017training}, a parameter-free coin betting optimizer that eliminates the need for learning-rate tuning. The coin-betting framework, developed by \citet{orabona2016coin}, reformulates online learning as a sequential betting problem and yields regret bounds that automatically adapt to unknown gradient scales, with subsequent work extending these guarantees to general scale-free and adaptive settings \citep{mcmahan2014unconstrained, cutkosky2018blackbox, orabona2018scalefree, jun2017changing, jun2017improved}. Despite these theoretical advances, coin-betting methods have not yet been applied to volatility forecasting or tail-risk calibration.

We evaluate MAOL using daily returns and realized variance constructed from 5-minute intraday data, following \citet{liu2015does}, for 10-year U.S. Treasury futures over 1983-2024, a market that spans multiple major volatility regimes. We instantiate the framework with four base models (RV-GARCH, HAR-RV, Log-GARCH, Log-HAR) and adopt GARCH as the primary baseline because it provides a transparent, interpretable representation of volatility dynamics, allowing the role of the online learning layer to be clearly identified. For distributional forecasts, we start from a standard Student-t return distribution and apply the PIT-based recalibration layer to correct remaining misspecification. The choice of Student-t mainly serves as a simple baseline to illustrate the gains from the proposed method. The empirical analysis compares the proposed method with standard volatility benchmarks, including full-sample GARCH, full-sample HAR-RV, and the industry-standard EWMA (RiskMetrics) model.

Empirically, a fully sequential MAOL estimator achieves variance forecast accuracy essentially identical to a full-sample HAR-RV oracle for U.S. Treasury futures (average per-step regret $+0.00345$ QLIKE units over 40 years) and remains close to that frontier for German Bunds, UK Gilts, and Japanese JGBs, while outperforming EWMA by a wide margin (Diebold-Mariano statistic $-34.5$, $p < 0.001$). PIT-based recalibration materially improves distributional fit: VaR forecasts pass all four standard backtests at the 5\% level, with an empirical violation rate of 4.82\% and a Dynamic Quantile (DQ) test $p$-value of 0.412.

Beyond forecast accuracy, the adaptive layers recover economically interpretable dynamics: the shock-response parameter rises in turbulent periods, persistence falls when faster adaptation is needed, and the variance-level multiplier rises precisely when the base recursion underpredicts volatility during stress episodes, providing a transparent diagnostic of structural model inadequacy. A sample-length study on non-overlapping 2-15 year windows confirms stable performance throughout, and cross-country validation on three additional bond markets using identical settings confirms generalizability.

Taken together, the results indicate that the residual dynamics left unexplained by classical volatility models are not merely noise, but a systematic source of predictive information that can be exploited through online learning. More broadly, the paper highlights the potential of combining econometric modeling with online learning theory, offering a principled path toward self-learning financial forecasting systems that remain robust in the face of unknown and evolving market dynamics.

The remainder of the paper is organized as follows. Section 2 presents the MAOL framework.
Section 3 states the main theoretical guarantees, with proofs deferred to the appendix. Sections 4 and 5 describe the empirical strategy and results. Section 6 concludes.